\subsection{Rancangan Mekanisme Pembangkitan dan Penerimaan Audit Event}
\label{subsec:rancangan-pembangkitan-event}

Berdasarkan analisis pada Subbab~\ref{sec:analisis-event-source} dan
analisis mekanisme pengiriman log pada
Subbab~\ref{subsubsec:analisis-mekanisme-pengiriman-log}, \textit{audit
event} dibangkitkan dari empat \textit{event source} yang telah
diidentifikasi, yaitu \textit{Patient Client}, \textit{Hospital Client},
\textit{Ministry Client}, dan PRE Server. Rancangan mekanisme ini mencakup
tiga aspek utama yang saling berkaitan: enkapsulasi \textit{event} dalam
format terenkripsi dan tertandatangani untuk menjamin kerahasiaan dan
\textit{non-repudiation}, mekanisme pengiriman dengan jaminan ketercapaian
untuk memastikan tidak ada \textit{event} yang hilang, serta penerimaan dan
verifikasi pada sisi ATS.

\subsubsection{Rancangan Enkapsulasi Audit Event}
\label{subsubsec:rancangan-enkapsulasi-event}

Analisis pada Subbab~\ref{subsubsec:analisis-jaminan-keasilan-log}
menetapkan bahwa setiap \textit{event source} harus membungkus
\textit{audit event} dalam format yang menjamin keaslian pengirim dan
\textit{non-repu- diation}. Namun, terdapat kebutuhan keamanan tambahan yang
perlu dipenuhi, yaitu kerahasiaan isi \textit{event} selama transit. \textit{Audit event} memuat atribut
yang berpotensi sensitif, seperti identitas aktor, jenis tindakan yang
dilakukan, dan informasi objek yang diakses. Meskipun kanal komunikasi
umumnya dilindungi oleh TLS, perlindungan pada lapisan transport bersifat
\textit{hop-by-hop} dan tidak memberikan jaminan kerahasiaan yang melekat
pada data itu sendiri setelah transit selesai.

Untuk memenuhi kedua kebutuhan tersebut secara bersamaan, rancangan
mengadopsi pendekatan \textit{Encrypt-then-Sign}, yaitu \textit{payload
event} dienkripsi terlebih dahulu kemudian \textit{ciphertext} yang
dihasilkan ditandatangani. Pola ini dipilih karena tanda tangan diberikan
atas \textit{ciphertext}, bukan \textit{plaintext}, sehingga secara
kriptografis membuktikan bahwa pengirim yang menciptakan \textit{ciphertext}
tersebut dan bertanggung jawab atas isinya, bukan sekadar memiliki akses
atas informasi yang terkandung di dalamnya. Rancangan enkapsulasi terdiri
atas dua lapisan mekanisme yang diuraikan pada subbab-subbab berikut.

\paragraph{Lapisan Enkripsi: AES-256-GCM dengan Distribusi Kunci ECIES}

\textit{Payload audit event} yang telah diserialisasi ke JSON dienkripsi
menggunakan algoritma AES-256-GCM dengan kunci sesi yang dibangkitkan secara
acak pada setiap pengiriman. AES-256-GCM dipilih karena bersifat
\textit{authenticated encryption}, yaitu selain menjamin kerahasiaan,
algoritma ini juga menghasilkan \textit{authentication tag} yang memastikan
integritas \textit{ciphertext} tidak dapat dirusak tanpa terdeteksi.

Distribusi kunci sesi AES ke ATS dilakukan menggunakan mekanisme ECIES
(\textit{Elliptic Curve Integrated Encryption Scheme}) berbasis kurva P-256.
Alur distribusi kunci dirancang sebagai berikut:

\begin{enumerate}
    \item Pengirim membangkitkan pasangan kunci \textit{ephemeral} P-256
    untuk setiap pengiriman.
    \item Pengirim melakukan operasi ECDH antara kunci privat \textit{ephemeral}
    dengan kunci publik ATS yang telah dikonfigurasi sebelumnya, menghasilkan
    \textit{shared secret}.
    \item \textit{Shared secret} diturunkan menjadi kunci pembungkus
    (\textit{wrapping key}) 256-bit menggunakan KDF SHA-256.
    \item Kunci sesi AES dienkripsi menggunakan kunci pembungkus tersebut
    dengan AES-256-GCM, menghasilkan kunci sesi terenkripsi.
    \item Kunci publik \textit{ephemeral}, \textit{nonce} pembungkus, dan
    kunci sesi terenkripsi digabungkan menjadi satu bidang \texttt{enc\_aes\_key}
    dalam format \texttt{<base64(epheme- ral\_pubkey)>.<base64(wrap\_nonce)>
    .<base64(enc\_aes\_key)>}.
\end{enumerate}

Dengan skema ini, kunci sesi AES hanya dapat dipulihkan oleh ATS menggunakan
kunci privatnya melalui operasi ECDH yang berkebalikan, tanpa memerlukan
pertukaran kunci awal (\textit{key exchange}) yang terpisah antara pengirim
dan penerima. Setiap pengiriman menggunakan kunci \textit{ephemeral} yang
berbeda, sehingga kompromi atas satu kunci sesi tidak mengancam kerahasiaan
\textit{event} lainnya (\textit{forward secrecy}).

\paragraph{Lapisan Penandatanganan: Ed25519 dengan Kunci Privat IOTA Aktor}

Setelah \textit{payload} dienkripsi, \textit{ciphertext} yang dihasilkan
ditandatangani menggunakan kunci privat IOTA Ed25519 dari aktor yang memicu
\textit{event}. Penggunaan kunci privat IOTA aktor, bukan kunci yang
dibangkitkan sementara oleh modul ATS pada komponen pengirim, merupakan
keputusan rancangan yang krusial untuk memperkuat sifat
\textit{non-repudiation}.

Dasar pertimbangannya adalah sebagai berikut. Kunci privat IOTA setiap aktor
telah terikat secara kriptografis dengan identitas aktor melalui mekanisme
BIP39 dan derivasi kunci yang digunakan DecMed; setiap aktor memiliki satu
pasang kunci yang konsisten sepanjang siklus hidupnya dalam sistem dan yang
sama digunakannya untuk menandatangani transaksi IOTA yang tercatat secara
\textit{on-chain}. Dengan demikian, tanda tangan pada \textit{audit event}
dapat dihubungkan langsung ke \textit{IOTA address} pengirim yang tercatat
pada \textit{ledger}. Apabila suatu aktor menyangkal telah membangkitkan
sebuah \textit{audit event}, tanda tangan yang ada pada \textit{event}
tersebut dapat diverifikasi terhadap kunci publik yang terhubung ke
\textit{IOTA address} yang sama yang digunakan aktor tersebut untuk
bertransaksi di jaringan IOTA, sehingga penyangkalan menjadi tidak dapat
diterima selama kunci privat tidak dikompromikan.

Tanda tangan Ed25519 dibuat menggunakan skema \textit{personal message} yang
konsisten dengan mekanisme penandatanganan pada ekosistem IOTA, yaitu pesan
yang ditandatangani adalah \textit{byte-byte ciphertext} yang dibungkus dalam
\texttt{IntentMessage} dengan \texttt{Intent::personal\_message()}. Hasilnya
diserialisasi ke dalam representasi Base64.

\paragraph{Format Akhir: EncryptedSignedEvent}

Hasil akhir enkapsulasi adalah objek \texttt{EncryptedSignedEvent} yang
memuat lima bidang sebagaimana ditunjukkan pada
Tabel~\ref{tab:encrypted-signed-event}.

\begin{table}[htbp]
\centering
\caption{Bidang pada Struktur \texttt{EncryptedSignedEvent}}
\label{tab:encrypted-signed-event}
\renewcommand{\arraystretch}{1.3}
\begin{tabular}{|p{3.5cm}|p{8.5cm}|}
\hline
\textbf{Bidang} & \textbf{Deskripsi} \\
\hline
\texttt{enc\_aes\_key} &
Kunci sesi AES yang dienkripsi dengan ECIES, dalam format
\texttt{<ephemeral\_pubkey>.<wrap\_nonce>.<enc\_key>} yang masing-masing
dikodekan Base64 dan dipisahkan titik. \\
\hline
\texttt{ciphertext} &
\textit{Payload audit event} yang dienkripsi menggunakan AES-256-GCM,
dikodekan Base64. \\
\hline
\texttt{nonce} &
\textit{Nonce} AES-GCM sepanjang 12 byte, dikodekan Base64. \\
\hline
\texttt{signature} &
Tanda tangan Ed25519 atas \textit{byte-byte ciphertext}, dikodekan Base64. \\
\hline
\texttt{iota\_address} &
\textit{IOTA address} pengirim yang terhubung ke kunci publik Ed25519 yang
digunakan untuk membuat tanda tangan, digunakan oleh ATS untuk keperluan
verifikasi. \\
\hline
\end{tabular}
\end{table}

Skema ini memisahkan secara tegas antara mekanisme kerahasiaan (ECIES +
AES-256-GCM) dan mekanisme \textit{non-repudiation} (Ed25519 atas
\textit{ciphertext}), sehingga kedua properti keamanan dapat diverifikasi
secara independen oleh pihak yang berwenang.

\subsubsection{Rancangan Mekanisme Pengiriman dengan Jaminan Ketercapaian}
\label{subsubsec:rancangan-pengiriman-event}

Analisis pada Subbab~\ref{subsubsec:analisis-mekanisme-pengiriman-log}
menetapkan bahwa pengiriman \textit{event} harus bersifat
\textit{asynchronous push-based} agar tidak mengganggu latensi layanan
utama. Namun, sifat asinkron menimbulkan risiko kehilangan \textit{event}
apabila ATS tidak dapat dijangkau pada saat \textit{event} dibangkitkan,
misalnya karena gangguan jaringan sementara atau ATS sedang tidak beroperasi.
Kegagalan pencatatan \textit{event} yang tidak terdeteksi merupakan ancaman
terhadap kelengkapan audit trail sebagaimana disyaratkan oleh KNF-TA-04,
sehingga mekanisme pengiriman harus memberikan jaminan ketercapaian
(\textit{delivery guarantee}) yang kuat di samping tetap bersifat
\textit{non-blocking}.

Untuk menjawab kebutuhan ini, rancangan mengadopsi pola
\textit{persistent queue} pada sisi pengirim yang digabungkan dengan
pengiriman asinkron. Perbandingan antara pendekatan pengiriman murni
\textit{fire-and-forget} dan pendekatan \textit{persistent queue} yang
diusulkan ditunjukkan pada Tabel~\ref{tab:analisis-pengiriman-event}.

\begin{table}[htbp]
\centering
\caption{Analisis Pendekatan Mekanisme Pengiriman \textit{Audit Event}}
\label{tab:analisis-pengiriman-event}
\renewcommand{\arraystretch}{1.3}
\begin{tabular}{|p{3cm}|p{4.5cm}|p{5cm}|}
\hline
\textbf{Pendekatan} & \textbf{Keunggulan} & \textbf{Keterbatasan} \\
\hline
\textit{Fire-and-Forget} Murni &
\textit{Non-blocking}, tidak menambah kompleksitas komponen pengirim. &
\textit{Event} hilang secara permanen apabila ATS tidak dapat dijangkau
saat pengiriman; tidak ada mekanisme pemulihan. \\
\hline
\textit{Persistent Queue} + Pengiriman Asinkron &
Menjamin setiap \textit{event} akan terkirim pada akhirnya melalui
mekanisme \textit{retry}; tahan terhadap gangguan jaringan sementara dan
\textit{restart} komponen. &
Membutuhkan pengelolaan berkas antrian lokal dan \textit{background worker}
pada setiap \textit{event source}. \\
\hline
\end{tabular}
\end{table}

Berdasarkan Tabel~\ref{tab:analisis-pengiriman-event}, pendekatan
\textit{persistent queue} dipilih karena jaminan ketercapaian yang
diberikannya jauh lebih penting bagi integritas audit trail dibandingkan
kompleksitas implementasi tambahan yang ditimbulkannya. Alur pengiriman
\textit{event} yang dirancang adalah sebagai berikut.

\begin{enumerate}
    \item \textbf{Konstruksi \texttt{EncryptedSignedEvent}.} Setiap kali
    sebuah \textit{event} dibangkitkan, komponen modul \texttt{ats} pada
    \textit{event source} mengonstruksi \texttt{EncryptedSignedEvent} sesuai
    mekanisme enkapsulasi pada
    Subbab~\ref{subsubsec:rancangan-enkapsulasi-event}.

    \item \textbf{Penyimpanan ke antrian persisten.} Sebelum percobaan
    pengiriman apa pun dilakukan, \texttt{EncryptedSignedEvent} yang telah
    dikonstruksi disimpan terlebih dahulu ke dalam berkas antrian persisten
    berbasis \textit{append-only} pada penyimpanan lokal \textit{event source}.
    Berkas antrian menggunakan format JSON-Lines, di mana setiap baris merupakan
    satu entri antrian (\texttt{QueueEntry}) yang memuat \textit{signed payload},
    pengenal unik entri, label konteks, waktu pembuatan, dan jumlah percobaan
    pengiriman yang telah dilakukan. Penyimpanan ke antrian dilakukan terlebih
    dahulu agar \textit{event} aman tersimpan secara lokal terlepas dari hasil
    pengiriman berikutnya.

    \item \textbf{Percobaan pengiriman langsung.} Setelah entri berhasil
    tersimpan di antrian, pengirim segera mencoba mengirimkan
    \texttt{EncryptedSignedEvent} ke \textit{endpoint}
    \texttt{POST /api/events} pada ATS secara langsung dalam \textit{task}
    asinkron yang tidak memblokir alur layanan utama. Apabila pengiriman
    langsung berhasil (ATS mengembalikan respons sukses), entri yang
    bersangkutan dihapus dari antrian.

    \item \textbf{Pengiriman ulang oleh \textit{background retry worker}.}
    Apabila pengiriman langsung gagal karena ATS tidak dapat dijangkau atau
    mengembalikan respons \textit{error}, entri tetap berada di antrian.
    \textit{Background retry worker} yang berjalan secara periodik pada setiap
    \textit{event source} membaca seluruh entri yang tersisa di antrian dan
    mencoba mengirimkannya kembali. \textit{Worker} ini menggunakan strategi
    \textit{exponential backoff} untuk menentukan interval antar-percobaan,
    sehingga menghindari pembanjiran ATS saat baru pulih dari gangguan.
    Entri yang telah melebihi batas percobaan pengiriman dikecualikan
    sementara dari percobaan ulang, tetapi tetap disimpan di antrian agar
    tidak hilang secara permanen dan dapat diinspeksi oleh operator.
\end{enumerate}

Rancangan ini menjamin bahwa setiap \textit{event} yang berhasil dikonstruksi
dan disimpan ke antrian akan terkirim ke ATS pada akhirnya, bahkan jika
\textit{event source} mengalami \textit{restart} di antara percobaan
pengiriman, selama berkas antrian tidak rusak secara fisik. Pemisahan antara
tahap penyimpanan ke antrian dan tahap pengiriman ke ATS juga memastikan bahwa
kegagalan jaringan tidak menyebabkan \textit{event} hilang secara diam-diam
tanpa sepengetahuan operator.

\subsubsection{Rancangan Penerimaan dan Verifikasi Event pada ATS}
\label{subsubsec:rancangan-penerimaan-event}

Pada sisi ATS, \texttt{EncryptedSignedEvent} yang diterima melalui
\textit{endpoint} \texttt{POST /api/events} diproses dalam dua fase sebelum
diteruskan ke komponen Audit Logger.

\paragraph{Fase Pertama: Verifikasi Tanda Tangan}

ATS mengekstrak \texttt{iota\_address} dari \textit{event} yang diterima,
kemudian memverifikasi tanda tangan Ed25519 yang terdapat pada bidang
\texttt{signature} terhadap \textit{byte-byte ciphertext} pada bidang
\texttt{ciphertext}. Verifikasi menggunakan kunci publik Ed25519 yang
diturunkan dari \texttt{iota\_address} pengirim. Apabila verifikasi
berhasil, ATS dapat menyatakan bahwa \textit{event} berasal dari pemegang
kunci privat yang terhubung ke \texttt{iota\_address} yang diklaim, dan bahwa
\textit{ciphertext} tidak mengalami perubahan sejak ditandatangani. Apabila
verifikasi gagal, \textit{event} ditolak tanpa diproses lebih lanjut,
mencegah injeksi \textit{event} palsu ke dalam rantai audit.

Pada fase ini, ATS tidak mendekripsi isi \textit{event}. Verifikasi tanda
tangan dilakukan sepenuhnya terhadap \textit{ciphertext}, konsisten dengan
pola \textit{Encrypt-then-Sign} yang diadopsi. Keputusan untuk memisahkan
fase verifikasi dari fase dekripsi merupakan keputusan rancangan yang disengaja:
verifikasi adalah hak ATS sebagai penjaga gerbang pencatatan, sedangkan
dekripsi adalah hak pihak yang berwenang melakukan audit atau investigasi.

\paragraph{Fase Kedua: Penerusan ke Antrian Internal}

\texttt{EncryptedSignedEvent} yang telah terverifikasi tanda tangannya
dikirimkan ke antrian asinkron internal yang menghubungkan komponen Audit
Receiver dengan komponen Audit Logger. Antrian ini diimplementasikan sebagai
kanal \textit{multi-producer single-consumer} (mpsc) dengan kapasitas yang
cukup untuk menyerap lonjakan kedatangan \textit{event}. Desain ini memisahkan
proses penerimaan HTTP dari proses pembentukan dan penulisan \textit{audit
record}, sehingga keduanya dapat berjalan secara konkuren tanpa saling
memblokir. Antrian internal bertindak sebagai penyangga (\textit{buffer})
yang mencegah lonjakan \textit{event} menghambat penerimaan \textit{event}
baru berikutnya.

\subsubsection{Rancangan Penyimpanan Audit Record dalam Kondisi Terenkripsi}
\label{subsubsec:rancangan-penyimpanan-terenkripsi}

Setelah \texttt{EncryptedSignedEvent} terverifikasi dan diteruskan oleh Audit
Logger, \textit{event} tersebut disimpan ke dalam \textit{audit record} dan
kemudian ke berkas log lokal serta IPFS dalam kondisi masih terenkripsi.
Keputusan rancangan untuk tidak mendekripsi \textit{event} pada titik ini
didasarkan pada dua pertimbangan berikut.

Pertama, pemisahan tanggung jawab. ATS berperan sebagai penyimpan bukti
digital yang andal, bukan sebagai pemroses isi \textit{event}. Mendekripsi
\textit{event} sebelum penyimpanan akan memperluas permukaan serangan karena
kunci privat ATS harus tersedia secara aktif setiap kali \textit{event}
diproses. Dengan menyimpan \textit{event} dalam kondisi terenkripsi, kunci
privat ATS hanya perlu tersedia saat audit atau investigasi dilakukan.

Kedua, kerahasiaan jangka panjang. Berkas log yang tersimpan di IPFS dapat
diakses oleh pihak mana pun yang mengetahui CID-nya. Apabila \textit{event}
disimpan dalam kondisi \textit{plaintext}, kerahasiaan isi audit trail
bergantung sepenuhnya pada keamanan akses ke IPFS. Dengan menyimpan
\textit{event} dalam kondisi terenkripsi, kerahasiaan isi audit trail
dijamin secara kriptografis: hanya pemegang kunci privat ATS yang dapat
mendekripsi dan membaca isi \textit{event}, terlepas dari siapa yang dapat
mengakses berkas log di IPFS.

Implikasi praktis dari keputusan ini adalah bahwa mekanisme
\textit{hash-chaining} pada Subbab~\ref{subsec:rancangan-pembentukan-record}
diterapkan terhadap \texttt{EncryptedSignedEvent} sebagai satu kesatuan,
bukan terhadap isi \textit{plaintext event}. Nilai \texttt{record\_hash}
dihitung dari kombinasi \texttt{record\_id}, \texttt{timestamp},
\texttt{prev\_record\_hash}, dan seluruh bidang \texttt{EncryptedSignedEvent}
yang mencakup \texttt{ciphertext} dan \texttt{signature}. Dengan demikian,
integritas rantai \textit{audit record} tetap dapat diverifikasi tanpa perlu
mendekripsi isi \textit{event} terlebih dahulu: perubahan pada
\textit{ciphertext} mana pun akan merusak \texttt{record\_hash} dan
mengakibatkan kegagalan verifikasi rantai, sebagaimana berlaku pada
mekanisme \textit{hash-chaining} yang telah dianalisis pada
Subbab~\ref{subsubsec:analisis-hash-chaining}.